% Part: applied-modal-logics
% Chapter: epistemic-logic
% Section: truth-at-w

\documentclass[../../../include/open-logic-section]{subfiles}

\begin{document}

\olfileid{aml}{el}{trw}

\olsection{صدق در یک جهان}

درست همانند منطق موجهات عادی، هر مدل معرفتی تعیین می‌کند کدام
!!{formula} در کدام جهان‌های آن صادق به شمار می‌آیند. برای مفهوم پایهٔ
معناشناسی رابطه‌ای از همان نمادگذاریِ «مدل~$\mModel{M}$،
!!{formula}~$!A$ را در جهان~$w$ صادق می‌سازد» استفاده می‌کنیم. این رابطه
به‌طور استقرایی تعریف می‌شود و برای همهٔ عملگرهای غیرموجهاتی با حالت
منطق موجهات عادی یکسان است.

\begin{defn}\ollabel{defn:mmodels}
  \emph{صدق !!a{formula}~$!A$ در~$w$} در مدل~$\mModel M$، که با
  $\mSat{M}{!A}[w]$ نمادگذاری می‌شود، به‌طور استقرایی چنین تعریف می‌شود:
  \begin{enumerate}
  \tagitem{prvFalse}{\indcase{!A}{\lfalse}{$\mSat{M}{\lfalse}[w]$ هرگز برقرار نیست.}}{}
  \tagitem{prvTrue}{\indcase{!A}{\ltrue}{$\mSat{M}{\ltrue}[w]$ همواره برقرار است.}}{}
  \item $\mSat{M}{p}[w]$ اگر و تنها اگر~$w \in V(p)$.
  \tagitem{prvNot}{\indcase{!A}{\lnot !B}{$\mSat{M}{\indfrm}[w]$ اگر و تنها اگر
    $\mSat/{M}{!B}[w]$.}}{}
  \tagitem{prvAnd}{\indcase{!A}{(!B \land !C)}{$\mSat{M}{\indfrm}[w]$ اگر و تنها اگر
    $\mSat{M}{!B}[w]$ و~$\mSat{M}{!C}[w]$.}}{}
  \tagitem{prvOr}{\indcase{!A}{(!B \lor !C)}{$\mSat{M}{\indfrm}[w]$ اگر و تنها اگر
    $\mSat{M}{!B}[w]$ یا~$\mSat{M}{!C}[w]$} (یا هر دو).}{}
  \tagitem{prvIf}{\indcase{!A}{(!B \lif !C)}{$\mSat{M}{\indfrm}[w]$ اگر و تنها اگر
    $\mSat/{M}{!B}[w]$ یا~$\mSat{M}{!C}[w]$.}}{}
  \tagitem{prvIff}{\indcase{!A}{(!B \liff !C)}{$\mSat{M}{\indfrm}[w]$ اگر و تنها اگر
    یا هر دو رابطهٔ~$\mSat{M}{!B}[w]$ و~$\mSat{M}{!C}[w]$ برقرار باشند، یا
    هیچ‌یک از~$\mSat{M}{!B}[w]$ و~$\mSat{M}{!C}[w]$ برقرار نباشند.}}{}
  \item\ollabel{defn:sub:mmodels-box}
    \indcase{!A}{\Knows_a !B}{$\mSat{M}{\indfrm}[w]$ اگر و تنها اگر
    برای هر~$w' \in W$ که~$R_a ww'$، رابطهٔ~$\mSat{M}{!B}[w']$ برقرار باشد.}
  \end{enumerate}
\end{defn}

بااین‌حال، اینجا باید دربارهٔ قیود روابط دسترسی‌پذیری خود بیندیشیم.
زیرا بنا بر بند~\olref{defn:sub:mmodels-box}، هرگاه هیچ~$w'$ای نباشد
که~$R_a ww'$، !!a{formula}~$\Knows_a !B$ در~$w$ صادق است. این همان بند
منطق موجهات عادی است: وقتی جهانی هیچ جانشینی ندارد، همهٔ فرمول‌های
$\Box$ در آن به‌طور تهی صادق‌اند. اما اگر~$\Knows$ را نمایانگر
\emph{دانش} بدانیم، این نتیجه به‌شدت خلاف شهود می‌نماید. به‌هرحال، معمولاً
می‌پنداریم هیچ وضعیتی وجود ندارد که در آن عاملی بتواند هم‌زمان هم~$!A$
و هم~$\lnot !A$ را بداند.

یک راه‌حل آن است که تضمین کنیم رابطهٔ دسترسی‌پذیری در منطق معرفتی همواره
\emph{بازتابی} باشد. این تقریباً متناظر با این اندیشه است که جهان واقعی
با اطلاعات عامل سازگار است. در واقع، منطق‌های معرفتی معمولاً از~S5
استفاده می‌کنند، اما بسته به اینکه دقیقاً بخواهند رابطهٔ~$\Knows_a$ چه
چیزی را نمایان سازد، ممکن است دستگاه‌های ضعیف‌تری به کار گیرند.

\begin{figure}
  \begin{center}
    \begin{tikzpicture}[modal]
      \node[world] (w1) [label={[align=right]right:\mTrue{p}\\ \mFalse{q}}]{$w_1$};
      \node[world] (w2) [label={[align=right]right:\mTrue{p}\\ \mTrue{q}},
        above right=of w1]{$w_2$};
      \node[world] (w3) [label={[align=right]right:\mFalse{p}\\ \mFalse{q}},
        below right=of w1] {$w_3$};
      \draw[<->] (w1) to node {$a$} (w2);
      \draw[->,loop left] (w1) to node {$a,b$} (w1);
      \draw[<->] (w1) to [swap] node {$b$} (w3);
      \draw[->,loop above] (w2) to node {$a,b$} (w2) ;
      \draw[->,loop below] (w3) to node {$a,b$} (w3) ;
    \end{tikzpicture}
  \end{center}
  \caption{یک مدل معرفتی ساده.}
  \ollabel{fig:simple}
\end{figure}

\begin{prob}
  بررسی کنید کدام‌یک از موارد زیر در
  \olref[aml][el][trw]{fig:simple} برقرارند:
  \begin{enumerate}
    \item $\mSat{M}{\lnot q}[w_1]$؛
    \item $\mSat{M}{\Knows_a \lnot q}[w_1]$؛
    \item $\mSat{M}{\Knows_b \lnot q}[w_1]$؛
    \item $\mSat{M}{\Knows_b q \lor \Knows_b \lnot q}[w_2]$؛
    \item $\mSat{M}{\Knows_a( \Knows_b q \lor \Knows_b \lnot q)}[w_2]$؛
    \item $\mSat{M}{\EKnows_{\{a,b\}} \lnot q}[w_3]$؛
    \end{enumerate}
\end{prob}

اکنون که تعریف پایهٔ صدق در یک جهان را داده‌ایم، دیگر مفاهیم معناشناختیِ
منطق موجهات عادی، مانند اعتبار موجهاتی و استلزام، بی‌درنگ به اینجا منتقل
می‌شوند و بر این شیوهٔ تازهٔ اندیشیدن به تفسیر عملگرهای موجهاتی اعمال
می‌گردند.

اکنون همچنین می‌توانیم شرایط صدق عملگر دانش مشترک~$\CKnows_G$ را بیان
کنیم. به یاد آورید که در~\olref[sfr][rel][ops]{sec}، \emph{بستار
تعدی}~$R^+$ رابطهٔ~$R$ چنین تعریف شد:
\begin{align*}
  R^+ &= \bigcup_{ n \in \mathbb{N}} R^n,
\intertext{که در آن}
  R^0 & = R \text{ و}\\
  R^{n+1} & = \Setabs{\tuple{x, z}}{\exists y (R^n xy \land Ryz)}.
\end{align*}
سپس، هرگاه~$G$ گروهی از عامل‌ها باشد،
$R_G = ( \bigcup_{b \in G} R_b )^+$ را بستار تعدیِ اجتماع روابط
دسترسی‌پذیری همهٔ عامل‌ها تعریف می‌کنیم.

\begin{defn}
اگر~$G' \subseteq G$، آنگاه~$\mSat{M}{\CKnows_{G'} !A}[ w]$ را اگر و
تنها اگر برای هر~$w'$ای که~$R_{G'} w w'$، رابطهٔ~$\mSat{M}{!A}[w']$
برقرار باشد، تعریف می‌کنیم.
\end{defn}

\end{document}
